\chapter*{Common Model and Notation}
\addcontentsline{toc}{chapter}{Common Model and Notation}

A single model is used throughout, and each chapter states any departure from it. Let $s\in S$ be an augmented economic state containing material resources, relevant institutional permissions, and outstanding obligations. In a state $s$ the available actions are $a\in A(s)$, and the allowed disturbances are $w\in W(s,a)$. Successor states are given by
\[
s' = T(s,a,w).
\]
Both $A(s)$ and $W(s,a)$ are assumed nonempty unless a result says otherwise. Let $G\subseteq S$ be the set of acceptable states.

A \emph{policy} is nonanticipating: it may use the observed history but cannot depend on disturbances that have not yet been observed. Unless stated, observation is complete, so the state and the past disturbances are observed. A policy keeps the system \emph{in $G$ forever} against a class of disturbances if every trajectory it generates under an allowed disturbance sequence stays in $G$.

Three structures are kept separate throughout: \emph{admissibility} (which trajectories satisfy the stated obligations), \emph{welfare} (which trajectories are better for whom), and \emph{financial value} (how trajectories are priced or accounted). A result about one is never silently a result about another.

\begin{center}
\small
\renewcommand{\arraystretch}{1.2}
\begin{tabular}{@{}ll@{}}
\toprule
Symbol & Meaning\\
\midrule
$S,\ A(s),\ W(s,a),\ T$ & states, actions, disturbances, transition map\\
$G$ & acceptable states\\
$\Phi$ & robust predecessor operator, $\Phi(X)=\{s\in G:\exists a\ \forall w,\ T(s,a,w)\in X\}$\\
$K=K(G)$ & viability kernel: greatest fixed point of $\Phi$\\
$L(h)$ & set of obligation-satisfying future sequences after history $h$\\
$h\sim h'$ & continuation equivalence, $L(h)=L(h')$\\
$F:S\to Y$ & an accounting or other representation\\
$\preceq$ & partial order on continuation profiles\\
\bottomrule
\end{tabular}
\end{center}
